% authority: science_draft
\documentclass[11pt]{article}

\usepackage[margin=1in]{geometry}
\usepackage{amsmath,amssymb,amsthm}
\usepackage{booktabs}
\usepackage{longtable}
\usepackage{enumitem}

\newcommand{\EqSrc}{\mathsf{Eq}_{\mathrm{src}}}
\newcommand{\EqRec}{\mathsf{Eq}_{\mathrm{src}}^{\mathrm{rec}}}
\newcommand{\EqIso}{\mathsf{Eq}_{\mathrm{src}}^{\mathrm{iso}}}
\newcommand{\Obj}{\mathsf{Obj}_{0}^{\mathrm{rec}}}
\newcommand{\Inv}{\mathsf{Inv}_{0}^{\mathrm{rec}}}
\newcommand{\Cmp}{\mathsf{Cmp}_{0}^{\mathrm{rec}}}
\newcommand{\Ann}{\mathsf{Ann}}
\newcommand{\Infl}{\mathsf{Infl}}
\newcommand{\Tok}{\mathsf{Tok}_{0}}
\newcommand{\Block}{\mathsf{Block}}
\newcommand{\Pass}{\mathsf{Pass}}
\newcommand{\RetainH}{\mathsf{Retain}_{H}}
\newcommand{\GenH}{\mathsf{Gen}_{H}}
\newcommand{\MetricData}{\mathrm{MetricData}(E)}
\newcommand{\geff}{g_{\mathrm{eff}}}

\newtheorem{definition}{Definition}
\newtheorem{theorem}{Theorem}
\newtheorem{lemma}{Lemma}
\newtheorem{obstruction}{Obstruction}
\newtheorem{boundary}{Boundary}
\newtheorem{nonconclusion}{Non-conclusion}

\title{Selected Upstream Equivalence Attempt v1}
\author{Ontology Formalizer Packet RT-20260706-011}
\date{2026-07-06}

\begin{document}
\maketitle

\section{Control Status}

This artifact implements v17 P6-T02. The P6-T01 selector chose
\(\EqSrc\_\mathrm{theorem\_attempt}\). The present packet executes that
attempt as a bounded theorem candidate over explicit source records and
source-supplied relabeling maps.

The result is not general \(\EqSrc\) discharge. It is not \(\RetainH\)
adoption, not \(\GenH\) adoption, not source-law adoption, not a canonical
ontology edit, not \(\MetricData\) adoption, not \(\geff\) scope expansion,
not matter coupling, not stress-energy semantics, not matter action, not
Einstein equations, not benchmark promotion, not a Gate Chair verdict, and
not completed derivation.

\section{Attempt Target}

\begin{definition}[Record-local EqSrc datum]
A record-local source-equivalence datum is a tuple
\[
  \mathcal E^{\mathrm{rec}}=(\Obj,\Inv,\Cmp,\Ann,\Infl)
\]
where \(\Obj\) is a finite set of declared source records, \(\Inv\) is a
finite source-only invariant ledger, \(\Cmp\) is a source-only comparison
rule, \(\Ann\) records comparison annotations, and \(\Infl\) records proxy
and forbidden-channel influence edges.
\end{definition}

\begin{definition}[Admissible relabeling map]
For \(X,Y\in\Obj\), an admissible relabeling map
\(\rho:X\to Y\) is source-side data satisfying all of the following:
\begin{enumerate}[leftmargin=*]
  \item \(\rho\) is bijective on declared source token identifiers.
  \item \(\rho\) preserves token class, annotation content, source path,
  object identifier, hash, rule scope, dependency, proxy-exclusion content,
  and influence-edge incidence.
  \item \(\rho\) preserves forbidden-channel orientation and proxy-class
  coverage.
  \item \(\rho\) does not use target topology, target atlas, target metric,
  proper time, detector protocol, stress-energy semantics, matter-action
  semantics, benchmark behavior, validation status, registry status, role
  identity, handoff state, or commit state as a premise.
  \item Missing annotation, unmapped proxy use, changed influence incidence,
  or a required target-side comparison returns \(\Block\).
\end{enumerate}
\end{definition}

\begin{definition}[Record-local equivalence relation candidate]
Write
\[
  X\EqRec Y
\]
when there exists an explicitly supplied admissible relabeling map
\(\rho:X\to Y\). The relation is defined only inside one declared finite
record orbit whose objects and relabeling maps are listed in the packet.
\end{definition}

\section{Theorem Candidate}

\begin{theorem}[Record-local EqSrc relabeling orbit theorem]
Let \(\mathcal O\subseteq\Obj\) be a finite declared source-record orbit.
Assume:
\begin{enumerate}[leftmargin=*]
  \item for every \(X\in\mathcal O\), the identity map \(1_X:X\to X\) is
  explicitly supplied and admissible;
  \item for every supplied admissible map \(\rho:X\to Y\), an explicitly
  supplied inverse map \(\rho^{-1}:Y\to X\) is admissible;
  \item for every pair of supplied admissible maps \(\rho:X\to Y\) and
  \(\sigma:Y\to Z\), an explicitly supplied composition
  \(\sigma\circ\rho:X\to Z\) is admissible;
  \item every identity, inverse, and composition check uses only the source
  invariants, annotations, and influence edges in
  \(\mathcal E^{\mathrm{rec}}\).
\end{enumerate}
Then \(\EqRec\) is an equivalence relation on \(\mathcal O\).
\end{theorem}

\begin{proof}
Reflexivity follows from premise 1: for every \(X\in\mathcal O\), the
explicit admissible identity \(1_X\) witnesses \(X\EqRec X\). Symmetry
follows from premise 2: if \(X\EqRec Y\), then an admissible \(\rho:X\to Y\)
is supplied, and the supplied admissible inverse \(\rho^{-1}:Y\to X\)
witnesses \(Y\EqRec X\). Transitivity follows from premise 3: if
\(X\EqRec Y\) and \(Y\EqRec Z\), supplied admissible maps
\(\rho:X\to Y\) and \(\sigma:Y\to Z\) exist, and the supplied admissible
composition \(\sigma\circ\rho:X\to Z\) witnesses \(X\EqRec Z\). Premise 4
keeps all three proof branches source-side. Therefore \(\EqRec\) is an
equivalence relation on the declared finite orbit only.
\end{proof}

\section{Finite Witness}

\begin{definition}[Two-object relabeling orbit]
Let \(\mathcal O=\{X_0,X_1\}\), where \(X_0\) is the manifest-bound
source-syntax object from the earlier local \(\EqIso\) datum and \(X_1\) is
its source-neutral relabeling. Let \(\rho:X_0\to X_1\) be the bijection that
renames source token identifiers while preserving token class, annotation
content, proxy coverage, forbidden-channel orientation, and influence-edge
incidence. Let \(\rho^{-1}\), \(1_{X_0}\), and \(1_{X_1}\) be explicitly
listed source maps.
\end{definition}

\begin{lemma}[Witness satisfies the theorem premises]
The finite orbit \(\{X_0,X_1\}\) satisfies the record-local theorem premises
when the supplied composition table is
\[
\begin{array}{c|c|c}
f & g & g\circ f \\
\hline
1_{X_0} & \rho & \rho \\
\rho & \rho^{-1} & 1_{X_0} \\
\rho^{-1} & \rho & 1_{X_1} \\
1_{X_1} & \rho^{-1} & \rho^{-1}
\end{array}
\]
with all rows checked only against source-token class, annotation, proxy, and
influence-edge invariants.
\end{lemma}

\begin{proof}
Each listed map is source-side data. The identity maps preserve every
source invariant by definition. The relabeling \(\rho\) and inverse
\(\rho^{-1}\) preserve the same source invariants because they are bijective
renamings of source token identifiers and do not alter annotation content or
influence incidence. The table supplies the only compositions needed inside
the two-object orbit, and every composition remains admissible by direct row
inspection. The theorem therefore applies.
\end{proof}

\section{Negative Controls}

If \(X_{\mathrm{ann}}\) weakens required annotation content, then
\[
  \Cmp(X_0,X_{\mathrm{ann}};\Inv)=\Block.
\]
If \(X_{\mathrm{proxy}}\) introduces an unmapped proxy or treats an omitted
influence edge as absent by default, then
\[
  \Cmp(X_0,X_{\mathrm{proxy}};\Inv)=\Block.
\]
These controls prevent the theorem candidate from becoming a notation
similarity test, target-success proxy, validation-status proxy, or
authority-label proxy.

\section{Scoped Obstruction To General EqSrc}

\begin{obstruction}[General EqSrc family closure is missing]
The record-local theorem does not prove general \(\EqSrc\). The obstruction
is
\[
  \mathrm{OB\mbox{-}P6T02\mbox{-}GENERAL\mbox{-}EQSRC\mbox{-}FAMILY\mbox{-}CLOSURE\mbox{-}MISSING}.
\]
A general theorem would need, for the relevant source families, a declared
family domain, an invariant ledger valid across the family, a source-only
comparison rule, explicit identity inverse and composition closure, negative
control retention, and proof that broadening does not require \(\RetainH\)
or \(\GenH\). Those objects are not supplied by the current packet.
\end{obstruction}

\section{RetainH and GenH Trigger Consequence}

This attempt keeps \(\RetainH\) and \(\GenH\) deferred for the fixed
record-local orbit. It exposes their conditional trigger boundary:
\(\RetainH\) becomes an immediate primitive only if a later route asks
whether source-record structure is preserved under an \(H\)-indexed retention
operation; \(\GenH\) becomes an immediate primitive only if a later route
constructs or ranges over an \(H\)-generated source family. Neither primitive
is adopted here.

\section{Distance-to-GR Effect}

\begin{boundary}
This packet supplies one bounded theorem candidate and one scoped obstruction.
It produces no Distance-to-GR ledger delta. The `source_equivalence_eqsrc'
ledger row remains draft/control with general equivalence theorem missing.
\end{boundary}

\begin{center}
\begin{tabular}{p{0.32\linewidth}p{0.56\linewidth}}
\toprule
\textbf{Burden} & \textbf{Status after P6-T02} \\
\midrule
Source ontology primitives & unchanged; no canonical ontology edit \\
Source equivalence EqSrc & record-local theorem candidate supplied; general EqSrc not discharged \\
RetainH & unchanged; not adopted; conditional trigger clarified \\
GenH & unchanged; not adopted; conditional trigger clarified \\
ObsLoc\_lc & unchanged \\
Resp\_lc & unchanged \\
M\_src & unchanged \\
g\_eff & unchanged; no scope expansion \\
matter coupling & unchanged; not derived and not adopted \\
Einstein equations & blocked \\
finite-variation robustness & unchanged \\
benchmark promotion & blocked \\
Gate Chair status & unchanged; no verdict issued \\
current route freeze or hard-fail status & not frozen; P6-T03 selector remains open \\
\bottomrule
\end{tabular}
\end{center}

\section{Next Route}

The logical next step is P6-T03: a theoretical-continuation-selector packet
should choose whether this attempt needs smuggling audit, Refuter stress,
repair, freeze review, or return-to-matter routing. The selector must not
overread the theorem candidate as general \(\EqSrc\).

\section{Forbidden Conclusions}

\begin{nonconclusion}
This theorem attempt does not authorize general \(\EqSrc\) discharge,
\(\RetainH\) adoption, \(\GenH\) adoption, source-law adoption,
\(\MetricData\) adoption, \(\geff\) scope expansion, physical metric
authority, detector calibration, matter semantics, detector semantics,
coupling-law adoption, matter-coupling derivation or adoption, stress-energy
semantics, stress-energy tensor, matter action, Einstein equations, benchmark
promotion, Gate Chair verdict, completed derivation, future source-extension
impossibility, or program-wide no-go conclusion.
\end{nonconclusion}

\section*{APA 7 Source Materials}

The AEther-Flow Research Project. (2026). \emph{Recommendations implementation
plan continue task v17} [Internal implementation plan].

The AEther-Flow Research Project. (2026). \emph{Upstream-burden selector v1}
[Internal science draft].

The AEther-Flow Research Project. (2026). \emph{Localization source-basis
axiom selector-domain EqSrc source-equivalence obligation packet or controlled
pause} [Internal science draft].

The AEther-Flow Research Project. (2026). \emph{Localization source-basis
axiom selector-domain EqSrc source-relabeling isomorphism datum} [Internal
science draft].

The AEther-Flow Research Project. (2026). \emph{V15 P5-T01 upstream dependency
audit for EqSrc RetainH and GenH} [Internal science draft].

\end{document}

